\section{NN-S008: a corrected spectral bridge and the scope of RH symmetry}
\label{sec:NN-S008}

\begin{proposition}[NN-S008]
The equality in corpus 09 and corpus 10, Theorem
\texttt{thm:bridge}, must be restricted to
\[
 \left\{z:\ |\operatorname{Im}z|<\tfrac12,
               \ \zeta(\tfrac12-iz)=0\right\},
\]
or written using $\xi$ in place of $\zeta$. Without the restriction
it includes trivial zeros and is false for the intended nontrivial
spectral support. Under A, the involutions $c(s)=\bar s$ and
$w(s)=1-s$ have identical actions on the reduced nontrivial zero
support. Their Klein-four action has image of order two and kernel
$\{1,cw\}$. This statement about support does not identify their
actions on the local jet algebras.
\end{proposition}

\begin{proof}
The trivial zero $\rho=-2$ maps under $z=i(\rho-1/2)$ to
$-5i/2$. It satisfies the unrestricted displayed zeta equation
and is outside the intended strip. Conversely the affine inverse
$s=1/2-iz$ sends the displayed strip exactly to $0<\Re s<1$,
so the restricted equality follows from the stated Connes--Consani
bridge. The completed entire $\xi$ already removes the trivial
zeros and poles. On the critical line $\bar\rho=1-\rho$, giving
$c=w$ on the support under A. Complex conjugation is nontrivial
there because there are nonreal zeros, so the image has order two.
At a fixed zero the composite still acts semilinearly on the
local coordinate by $z\mapsto-\bar z$; NN-S006 describes that
action on jets. The reduced-point calculation cannot erase it.
\end{proof}

\paragraph{Spectral and geometric dictionary.}
The Connes--Consani input is \emph{Hochschild homology, trace map
and $\zeta$-cycles}, equation (2.2), the common-zero assertion in
Section 2, and Proposition 4.2. Its spectral values become
$-\gamma^2-1/4$ under A. Real negative eigenvalues alone do not
prove self-adjointness or semisimplicity: a two-by-two Jordan
block with a real negative eigenvalue is a counterexample.
In geometric language the reduced support loses infinitesimal
and equivariant data; the real Artin schemes of NN-S006 retain
the latter. The passage from a scheme to its reduction is EGA I,
Section 5.1; no finite-field realization or canonical-basis
positivity follows from this reduction.

\section{NN-S009: a supported zero does not retain cancellation history}
\label{sec:NN-S009}

\begin{theorem}[NN-S009]
For any nonzero commutative ring $R$, the value in $G(R)$ of a
formal supported expression cannot distinguish a zero supplied
as a constant from a zero obtained by cancellation. The same
failure holds for a representation ring $R(\Gamma)$, even with
the grade or other ambient coordinate specified. Thus
``no supported zero'' is not an intrinsic equivalent of a
positive formula unless the encoding of absence and the allowed
formulas are specified.
\end{theorem}

\begin{proof}
The supported constant $e=0_R$ and the expression
$1\mathbin\oplus(-1)$ have the same value $e$ and the same
support bit $1$. In a representation ring the zero object and
the virtual expression $[V]-[V]$ likewise have the identical
supported value. More generally a class $r$ equals
$(r+a)-a$, so value and support cannot reveal whether an
expression used subtraction. Prescribing a grade does not change
any of these equalities; it only records that the value is to be
placed at that grade. Therefore any predicate of the resulting
$G(R)$-value takes the same answer on these pairs.

For comparison, with a \emph{specified witness encoding} the sum
$\bigoplus_{y\in W}1$ equals $\tau$ exactly when $W$ is empty.
Indeed a nonempty finite sum of supported terms remains supported,
even if its amplitude vanishes. This is a statement about that
encoding, and does not give a method of computing a sum over an
implicitly defined witness set.
\end{proof}

\paragraph{Application to GCT and moonshine.}
Corpus 11, its introductory four-way distinction, the two-zero-fibre
lemma, and the support-complete character inversion theorem correctly
separate a missing coordinate from a chosen zero coordinate.
They do not separate the last two constructions in the proof.
In particular a PH0 claim concerns effectivity of structure
constants in a chosen canonical basis or positive formula.
It is not a consequence of calling a coefficient supported.
PH1 requires further controlled counting geometry. The theorem
does not refute either correctly specified assertion.

\paragraph{Geometric dictionary for NN-S009.}
In the Grothendieck group of representations or vector bundles,
$[V]-[V]=0=[0]$ by its defining relations. The effective cone is
the image of actual objects and does not record a chosen
presentation of a virtual class. SGA 6, Expos\'e IV, Definition
equation (1.5.2) and Definition 2.2, provide this exact group-completion
dictionary. The support annotation is extra grading or
stratification data; it does not restore a lost exact-sequence
or expression history. The conclusion concerns both orbit-ring
multiplicities and representation-valued moonshine coefficients.

\section{NN-S010: the universal joint quotient for support and residue}
\label{sec:NN-S010}

\begin{theorem}[NN-S010]
For every integer $m\geq2$, the map
\[
 n\longmapsto(n\bmod m,\mathbf1_{n>0})
 \quad\text{from }\mathbb N\text{ to }(\mathbb Z/m\mathbb Z)\times\mathbb B
\]
has image isomorphic to $G(\mathbb Z/m\mathbb Z)$, with $m+1$
elements. It is the coarsest semiring quotient retaining both
residue and support: its congruence identifies precisely the
positive integers congruent modulo $m$, and leaves zero alone.
Neither residue alone nor support alone determines the other.
For $m=2$, addition is exactly the pair (XOR, OR).
\end{theorem}

\begin{proof}
Both coordinates are semiring homomorphisms: for nonnegative
integers positivity of a sum is disjunction and positivity of a
product is conjunction. Their joint image consists of $(0,0)$
and all $(a,1)$, since every residue has a positive representative.
The induced addition and multiplication are the pair normal form
of $G(\mathbb Z/m\mathbb Z)$. The kernel equivalence has precisely
the claimed classes. If another quotient retains both maps, two
integers identified in it must have identical joint images, so
it factors onto this image. This proves the universal claim.
Zero and $m$ have equal residue and different support; $1$ and
$m$ have equal support and different residues. In characteristic
two the residue sum is XOR and the second coordinate is OR.
\end{proof}

\paragraph{Trace and complexity consequence.}
A count divisible by $p$ is either zero or a positive multiple
of $p$. A mod-$p$ trace distinguishes neither case until support
is independently available. Conversely Boolean existence does
not compute a residue. For the natural one element there is
no unital semiring map $\mathbb B\to\mathbb F_p$, since
$1+1=1$ would force $1=0$; there is no unital map in the other
direction since summing $p$ copies of $1$ would send $0$ to $1$.
E supplies exact counts for its stated instance classes and
therefore both coordinates within its stated running time.
That time is subexponential, not polynomial; A supplies no
further general count-recovery procedure from these two maps.

\paragraph{Geometric dictionary for NN-S010.}
For a finite reduced scheme of $n$ rational points, the two
statistics are nonemptiness and the trace $n\bmod p$ on its
degree-zero coefficient space. The theorem identifies the
smallest semiring of joint values. The finite-fibre trace is
the sum over fibres in SGA 4, Expos\'e XVII, Theorem 6.2.3;
the coherent mod-$p$ point-count relation used elsewhere in
this catalogue is SGA 7 II, Expos\'e XXII, Theorem 3.1 and
Corollary 3.2. The abstract quotient itself neither gives a
polynomial-size cohomology presentation nor computes its trace.

\section{NN-S011: split determinant support is efficiently evaluable but not multiplicative}
\label{sec:NN-S011}

\paragraph{Inputs.}
Corpus 16, Theorem \texttt{thm:split-determinant}, identifies the
amplitude with the determinant and the support with the Boolean
permanent. The following computation and counterexample isolate
what that pair does and does not imply for prompt 12.

\begin{theorem}[NN-S011]
For a matrix with entries $\tau$ or $1$ in $G(\mathbb F_2)$,
its split determinant can be computed in polynomial time: its
amplitude is the parity of perfect matchings and its support is
their existence. The split determinant on all matrices is not
multiplicative, even in dimension two.
\end{theorem}

\begin{proof}
Compute the determinant of the amplitude matrix over $\mathbb F_2$
by Gaussian elimination. In this characteristic the determinant
is the permanent modulo two. Compute whether the bipartite graph
of supported entries has a perfect matching by the augmenting-path
algorithm. This algorithm starts at the empty matching and finds
an augmenting path by graph search, increasing cardinality each
time. If none exists, the symmetric difference with any putative
larger matching would contain such a path, so the current matching
is maximum. At most the matrix dimension many searches suffice.
The pair of answers is one of $(0,0),(0,1),(1,1)$ and determines
the split determinant. These are polynomially many field and
graph operations.

For nonmultiplicativity take
\[
 A=\begin{pmatrix}1&\tau\\1&\tau\end{pmatrix},\qquad
 B=\begin{pmatrix}1&1\\\tau&\tau\end{pmatrix}.
\]
Both split determinants are $\tau$, because both graphs lack a
perfect matching. Their matrix product has all four entries $1$.
Its two supported determinant terms sum to $e$, the supported
zero in characteristic two. Hence
$\det_{\rm sp}(AB)=e\ne\tau=
\det_{\rm sp}(A)\det_{\rm sp}(B)$.
\end{proof}

\paragraph{Geometric dictionary and circuit consequence.}
The ordinary determinant is the action on the top exterior power
of a free module; its multiplicativity is an identity in that
additive category. The Boolean permanent records a different
incidence relation, and the counterexample proves that this
support refinement is not a determinant functor with the same
composition axiom. The exterior-power operations are SGA 6,
Expos\'e V, Section 2. The specialized evaluation algorithm above
does not give a small monotone arithmetic circuit for the
universal support polynomial; it uses a graph algorithm. Thus
monotone circuit lower bounds for matching are compatible with
efficient evaluation of this supported-zero pair. No analogous
algorithm for arbitrary sparse quadratic satisfiability follows.
